\subsection{محدودیت‌ها و تحلیل شکست}
تحلیل عملکردی و معماری ماژول \lr{kmeans\_quantize} نقاط ضعف زیر را آشکار می‌سازد:
\begin{itemize}
    \item \textbf{سربار حافظه در مرحله بازیابی و جدول مراکز:} اگرچه ذخیره شاخص‌ها در ابعاد بیتی پایین صرفه‌جویی تئوریک ایجاد می‌کند، اما بازسازی ماتریس شناور در حین گذر پیشرو یا نگهداری بافرهای فعال‌سازی در گراف محاسباتی \lr{STE} نیازمند مصرف حافظه موقت بالا در کارت گرافیک است (\lr{2494.40 MB} در مدل \lr{DistilBERT}).
    \item \textbf{عدم کاهش زمان ضرب ماتریسی در غیاب هسته‌های اختصاصی:} در صورتی که سخت‌افزار از فرمت شاخص‌بندی‌شده به صورت سخت‌افزاری پشتیبانی نکند، عمل دکوانتایز برخط پیش از ضرب ماتریسی، مزیتی در شتاب‌بخشی محاسباتی (\lr{GEMM Speedup}) نسبت به محاسبات چگال شناور استاندارد ایجاد نخواهد کرد.
    \item \textbf{حساسیت بالا به انحراف توزیع مقادیر دورافتاده (\lr{Outliers}):} در لایه‌هایی که دارای مقادیر دورافتاده با دامنه بزرگ هستند، خوشه‌بندی یک‌بعدی مراکز زیادی را به پوشش کران‌های دور اختصاص داده و وضوح بازنمایی مقادیر حول میانگین صفر دچار افت می‌گردد.
\end{itemize}
